\section{Vấn đề nghiên cứu}
\label{sec:problem}

Các phương pháp tính toán để dự đoán lối sống thực khuẩn thể phải đối mặt với bốn thách thức chính. Thứ nhất, cơ sở dữ liệu hệ gen phage hiện có vẫn còn rất hạn chế so với sự đa dạng khổng lồ của phage trong tự nhiên, phần lớn chưa được phân lập và nghiên cứu đặc tính~\cite{hayes2017metagenomic}. Thứ hai, không giống vi khuẩn có gen 16S rRNA làm gen đánh dấu phổ quát, thực khuẩn thể không có gen đánh dấu tương tự~\cite{mokili2012metagenomics}, khiến các phương pháp phân loại dựa trên gen đánh dấu không khả thi. Thứ ba, dữ liệu metagenomic thường chứa nhiều đoạn DNA ngắn bao gồm chỉ một phần gen hoặc các chuỗi gen không hoàn chỉnh, làm phức tạp các chiến lược phân loại dựa trên gen chức năng~\cite{jiang2017ensvmb}. Thứ tư, thực khuẩn thể có thể chia sẻ các chuỗi gen với vi khuẩn chủ thông qua chuyển gen ngang, đặc biệt trong trường hợp prophage được tích hợp vào hệ gen vi khuẩn~\cite{rozov2017recycler}, dẫn đến phân loại sai trong quá trình gán phân loại học.

Dựa trên các thách thức trên, bài toán dự đoán lối sống thực khuẩn thể được phát biểu như sau: cho một tập hợp các contig $\mathcal{S} = \{S_1, S_2, \ldots, S_n\}$ thu được từ giải trình tự và lắp ráp metagenomic, trong đó mỗi contig $S_i$ có độ dài $L_i$ base pairs, nhiệm vụ là học hàm phân loại $f: \mathcal{S} \rightarrow \{0, 1\}$ trong đó $0$ đại diện cho lớp temperate và $1$ đại diện cho lớp virulent, sao cho $f$ hoạt động hiệu quả trên các contig có độ dài biến thiên lớn mà không phụ thuộc vào các công cụ dự đoán gen bên ngoài.

Các phương pháp hiện có đã đạt được tiến bộ đáng kể nhưng vẫn còn những hạn chế quan trọng. Các phương pháp học máy truyền thống như PHACTS~\cite{mcnair2012phacts}, PhageAI~\cite{tynecki2020phageai}, và BACPHLIP~\cite{hockenberry2021bacphlip} được thiết kế cho hệ gen hoàn chỉnh và hiệu suất giảm mạnh trên contig phân mảnh. DeePhage~\cite{wu2021deephage} sử dụng CNN trên mã hóa one-hot nhưng thiếu tri thức tiền huấn luyện, hạn chế khả năng tổng quát hóa trên các contig ngắn với thông tin ngữ cảnh hạn chế. PhaTYP~\cite{shang2023phatyp} cải thiện hiệu suất nhờ kiến trúc BERT nhưng phụ thuộc vào Prodigal~\cite{hyatt2010prodigal} để dự đoán protein—công cụ này có hiệu suất giảm đáng kể trên các đoạn ngắn hơn 200 bp~\cite{trimble2012short}, khiến PhaTYP kém hiệu quả trên contig ngắn thiếu vùng mã hóa protein hoàn chỉnh. DeepPL~\cite{zhang2024deeppl} và ProkBERT~\cite{ligeti2024prokbert} sử dụng tokenization k-mer chồng lấp, gây ra rò rỉ thông tin trong quá trình masked language modeling và kém hiệu quả tính toán do độ dài chuỗi token gần bằng độ dài chuỗi gốc~\cite{zhou2023dnabert}.
